\begin{table*}[h]
  \begin{tabularx}{\linewidth}{l*{3}{Y}*{3}{Y}}
\hline\hline
\vspace{0.1cm}
    Hadron species &  \multicolumn{3}{c}{\kstar} & \multicolumn{3}{c}{$\phi$} \\
    \pt (\GeVc)       &   0.4 & 3.0 & 10.0 &   0.6 & 3.0 & 10.0  \\
% Not obvious: would we prefer quoting only 2 and not 3 pt ranges here? 
% --- it might depend on the pT dendence of these systematics 
 %   PID technique  &   TPC-TOF & No PID & \multicolumn{2}{c}{TPC-TOF} \\
    \cmidrule(lr){2-4} \cmidrule(lr){5-7} 
    Global tracking efficiency & \multicolumn{3}{c}{6} & \multicolumn{3}{c}{6} \\
    Signal extraction &5.1&   4.6&   9.7&   3.1&   3.2&   8.5 \\ 
    Track selection cuts &    3.0 &  2.1 & 0.9 &  1.6 &   1.6 &  2.4  \\
    Particle identification & 1.8 & 2.5 &  4.0 &   1.1  & 1.9  &  2.1 \\ 
    Material budget     & 4.3 &  0.8 &  0.1 &  6.2 &   0.4 &    - \\ 
    Hadronic interactions & 1.9 &0.9 & 0.1 & 1.4 & 0.7 & - \\
    \cmidrule(lr){2-4} \cmidrule(lr){5-7} 
    Total & 9.8&  8.3&  12.1&  9.5& 7.3&    9.0 \\ 
    Total ($N_{ch}$-independent) & 7.7&6.6&8.1 & 3& 5&5 \\ 
\hline\hline
  \end{tabularx}
  \newline
    \caption{Main sources and values of the relative systematic uncertainties (expressed in \%) of the \pt-differential yields of $\phi$ and \kstar\ resonances. The values are reported for low, intermediate and high \pt. The The contributions that act differently in the various event classes are removed from the total (quadratic sum of all contributions), defining the $N_{ch}$-independent ones, which are correlated across different multiplicity intervals. }
  \label{tab:sysresonances}  
\end{table*}
